\scanpage{7}\mspage{6}
Thus, the shadow of a figure is formed of the loci which refine it -- one also says that these are the \uline{loci ``of the figure}'' (not to be confused, with the strata of the figure which are loci). The \unsure{\deleted{sh}multishadow|multishadow} is therefore \unsure{formed|formed} of the shadows of its strata, it is indexed by the unfolding $\widetilde F$, set of strata of $F$, the relation of \deleted{\unsure{of order|of inclusion}}\inserted{\unsure{inclusion}} of the set-theoretic figure being none other than the relation of \unsure{incidences|incidence} between strata of $F$, ie. the relation \unsure{of|in} $\widetilde F$ induced by $\leq$ \unsure{of|in} $\mathcal M$ (or by $\ll$ at will, it comes to the $\overline m$). Whereas a figure is practically never reconstructed starting from its shadow, it is nevertheless \deleted{\unsure{at|d}}\inserted{o}ften reconstructed (case of the ``set-theoretic \unsure{\deleted{w}W|W}orkshops'') \deleted{from}\inserted{starting} from its multishadow. We shall return to this.

\S{} There is reason to speak of the canonical order relations
